\section{The coordinate structure of the forest}\label{sec:forest}

We first establish the structural facts needed to separate a vector
into finitely many coordinates and its complement. The finite heights
give reflexivity and coordinate density; the ancestral order gives
contractive projections and an exact norm formula on the complementary
subtrees.

For a finite set $H$, the coordinate projections are
$P_Hu=\sum_{v\in H}u(v)e_v$ on $J$ and
$P_Hx=\sum_{v\in H}x(e_v)e_v^*$ on $X$; write $Q_H=I-P_H$.

\begin{lemma}\label{lem:forest-structure}
The space $X$ defined by \eqref{eq:segment-norm} is infinite dimensional,
separable, and reflexive, and the coordinate functionals have dense span
in $X$. Let $H\subset\mathcal F$ be finite and \emph{ancestral}, meaning
that it contains every ancestor of each of its vertices. The coordinate
projections $P_H$ and $Q_H$ are contractive on both $J$ and $X$.
Moreover, $P_Hx\to x$ along any increasing ancestral exhaustion of
$\mathcal F$.
\end{lemma}

\begin{proof}
Singleton tests give $\|u\|_2\le\rho(u)$. On a component of height
$h$, each segment has at most $h+1$ vertices. Thus Cauchy--Schwarz and
disjointness give
\[
 \|u\|_2\le\rho(u)\le\sqrt{h+1}\|u\|_2.
\]
Write $J_h$ for the completion on that component. Segment families
from different components combine freely, so
\[
 J=\left(\bigoplus_{h\ge1}J_h\right)_{\ell_2}.
\]
Each $J_h$ is isomorphic to Hilbert space. The dual of an $\ell_2$
sum is isometrically the $\ell_2$ sum of the duals: Cauchy--Schwarz
proves one inequality, and testing on almost norming unit vectors in
finitely many components, with proportional weights, proves the other.
Applying this identity twice proves reflexivity. One may also pass
from reflexivity of $J$ to that of its dual through the canonical maps.
Write $\iota_J:J\to J^{**}$ and $\iota_X:X\to X^{**}$.
For $\Phi\in X^{**}$, define $x=\Phi\circ\iota_J\in J^*=X$.
Every $f\in X^*=J^{**}$ is $\iota_J(p)$ for some $p\in J$, so
\[
 (\iota_Xx)(f)=f(x)=x(p)=\Phi(\iota_J(p))=\Phi(f).
\]
Thus $\iota_X$ is onto. In particular $X^*=J$ canonically. If the coordinate functionals did not have dense span in
$X$, Hahn--Banach would give a nonzero element of $X^*=J$ annihilating
them. All its coordinates would vanish, a contradiction. Countability
of $\mathcal F$ now proves separability; the independent coordinates
prove infinite dimensionality.

A segment meets $H$ in a prefix and its complement in a suffix.
Restricting disjoint segment tests to either set therefore proves
contractivity on $J$. The adjoints give precisely the coordinate
projections on $X$, and remain contractive. Finally, approximate $x$
by a finite coordinate combination. Once $H$ contains its support,
contractivity bounds the discarded tail of $x$ by that approximation
error. This proves convergence along an ancestral exhaustion.
\end{proof}

Let $G_H$ be the set of minimal vertices outside $H$. We call these
vertices \emph{gates}. They include roots of components disjoint from
$H$. Distinct gates are incomparable, and the descendant sets
$U_g=\{v:g\preceq v\}$ partition $\mathcal F\setminus H$.
Each segment crossing from $H$ to its complement passes through one
gate. Let $J(U_g)$ denote the coordinate subspace of $J$ supported in
$U_g$, with its inherited norm.

\begin{lemma}\label{lem:cones}
For every finite ancestral $H$,
\begin{equation}\label{eq:cone-decomposition}
 Q_HJ=\left(\bigoplus_{g\in G_H}J(U_g)\right)_{\ell_2},
 \qquad
 \|Q_Hx\|^2=\sum_{g\in G_H}\|x|_{J(U_g)}\|^2\quad(x\in X).
\end{equation}
The gate-coordinate map $x\mapsto(x(e_g))_{g\in G_H}$ is a contraction
from $X$ to $\ell_2(G_H)$.
\end{lemma}

\begin{proof}
On a vector supported outside $H$, every segment test reduces to a
segment in one cone. Grouping a disjoint family by its gates gives the
upper estimate for the first identity; combining almost maximizing
families from finitely many distinct cones gives the reverse estimate.
Pass to the completion. A functional supported outside $H$ has exactly
the norm of its restriction to $Q_HJ$, because $Q_H$ is contractive.
Duality for the $\ell_2$ sum therefore proves the second identity.

A segment contains at most one gate, so $\rho(u)=\|u\|_2$ for
vectors supported on $G_H$. Testing $x$ on such finite vectors proves
the gate-coordinate contraction.
\end{proof}
